% =====================================================================
\phantomsection
\addcontentsline{toc}{part}{Foundations}
\BookChapter[What Is a Programming Language?]{What Is a Programming Language?}
\label{ch:intro}

\begin{researchquestion}
What is a programming language, and why do engineers need so many different ones?
\end{researchquestion}
% =====================================================================

A computer is a machine that follows instructions very fast and very exactly. It does
not guess what you meant. A \term{programming language} is an artificial language,
with strict rules of spelling and grammar, that lets a person write those instructions
down in a form that can be turned into something the machine can execute. The written
instructions are called a \term{program}, or \term{source code}.

This chapter introduces the vocabulary we need to talk about languages: algorithms,
levels of abstraction, programming paradigms, type systems, and the main language
families. It ends with a short history and a table of which languages are used in which
branch of engineering.

\LearningOutcomes{Describe an algorithm separately from the language it is written in; place a language on the scale from machine code to domain-specific tools; recognise the main programming paradigms in a short piece of code; explain the difference between static and dynamic typing and between compiling and interpreting; and read a popularity ranking with the right amount of caution.}

\section{Algorithms come first}

Before there is a program there is an \term{algorithm}: a finite list of precise steps
that solves a problem. The idea is much older than computers. The Euclidean algorithm for
the greatest common divisor of two integers is more than 2000 years old:

\begin{enumerate}
\item Take two positive integers $a$ and $b$.
\item If $b = 0$, the answer is $a$. Stop.
\item Otherwise replace $(a, b)$ by $(b, a \bmod b)$, where $a\bmod b$ is the remainder
  of $a/b$, and go back to step 2.
\end{enumerate}

For $a = 1071$, $b = 462$ the steps give $(1071, 462) \to (462, 147) \to (147, 21)
\to (21, 0)$, so the answer is $21$. The algorithm is the same in every language. Here
it is in \py,
\begin{lstlisting}[language=Python]
def gcd(a, b):
    while b != 0:
        a, b = b, a % b
    return a
\end{lstlisting}
and here it is in C:
\begin{lstlisting}[language=C]
int gcd(int a, int b) {
    while (b != 0) {
        int r = a % b;
        a = b;
        b = r;
    }
    return a;
}
\end{lstlisting}
The two programs look different: C wants to know that \code{a} and \code{b} are
integers (\code{int}), it uses braces \code{\{ \}} to mark blocks and semicolons to end
statements, and it needs a temporary variable \code{r} for the swap. \py marks blocks
by indentation and can swap two values in one line. But the \emph{algorithm} is
identical. Keep this in mind when you read the rest of these notes: most of what
differs between languages is \emph{notation}; the ideas are shared.

\begin{keyidea}
Think about the algorithm first, and the language second. If you can describe the steps
of a solution clearly in words or with a flowchart, writing them in any language is a
matter of looking up the notation.
\end{keyidea}

\section{From problem to program: pseudocode and flowcharts}

Experienced programmers rarely start by typing code. They first describe the algorithm in
a form that is independent of any language. Two tools are common.

\paragraph{Pseudocode} is a precise, structured description in plain language, using the
building blocks that every language has (variables, decisions, loops). Here is Newton's
method for solving $f(x) = 0$, which we will meet in every language chapter:
\begin{lstlisting}[language={},basicstyle=\ttfamily\small,keywordstyle={},
  morekeywords={INPUT,OUTPUT,SET,FOR,TO,IF,THEN,RETURN,END,STOP,REPORT}]
INPUT  function f, derivative df, starting guess x0, tolerance tol, max_iter
SET x = x0
FOR k = 1 TO max_iter
    SET step = f(x) / df(x)
    SET x = x - step
    IF |step| < tol THEN
        RETURN x and k
    END IF
END FOR
REPORT "no convergence" and STOP
\end{lstlisting}

\paragraph{Flowcharts} show the same logic as a diagram (Figure~\ref{fig:flowchart}).
Ovals mark the start and end, rectangles are actions, diamonds are decisions, and arrows
show the order. Flowcharts are particularly helpful for logic with many branches, such as
the control sequence of a machine, and they are the ancestors of the graphical languages
described in Chapter~\ref{ch:others}.

\begin{figure}[htbp]
\centering
\begin{tikzpicture}[node distance=6mm and 10mm,
  every node/.style={font=\sffamily\small},
  term/.style={draw=nublue,fill=nublue!8,rounded corners=8pt,minimum width=26mm,minimum height=7mm},
  act/.style={draw=nublue,fill=white,minimum width=34mm,minimum height=7mm,align=center},
  dec/.style={draw=nugold!85!black,fill=nugold!10,diamond,aspect=2.4,inner sep=1pt,align=center},
  arr/.style={-{Stealth[length=2.2mm]},thick}]
\node[term] (start) {start};
\node[act,below=of start] (init) {read $f$, $f'$, $x_0$\\set $x = x_0$, $k = 0$};
\node[act,below=of init] (step) {$k = k + 1$\\$s = f(x)/f'(x)$, \ $x = x - s$};
\node[dec,below=of step] (conv) {$|s| < \mathrm{tol}$?};
\node[term,right=18mm of conv] (done) {output $x$, $k$};
\node[dec,below=of conv] (maxit) {$k \ge k_\mathrm{max}$?};
\node[term,right=18mm of maxit] (fail) {report failure};
\draw[arr] (start) -- (init);
\draw[arr] (init) -- (step);
\draw[arr] (step) -- (conv);
\draw[arr] (conv) -- node[above]{yes} (done);
\draw[arr] (conv) -- node[right]{no} (maxit);
\draw[arr] (maxit) -- node[above]{yes} (fail);
\draw[arr] (maxit.west) -- node[above]{no} ++(-14mm,0) |- (step.west);
\end{tikzpicture}
\caption{Flowchart of Newton's method with a stopping tolerance and a maximum number of
iterations.}
\label{fig:flowchart}
\end{figure}

Whichever tool you use, a good algorithm description answers four questions before any code
is written: What are the inputs, and in which units? What are the outputs? What steps
transform one into the other? And how will I check that the answer is right? The last
question is the one beginners most often forget.

\section{From switches to Python: levels of abstraction}

A processor (CPU) understands only \term{machine code}: sequences of binary numbers,
each of which encodes a very simple operation such as ``add the numbers in two
registers'' or ``jump to address 4096 if the last result was zero''. Writing programs
directly in machine code is possible but extremely tedious and error-prone. Over the
decades, programmers built layers of languages on top of machine code, each layer
hiding more detail.

\begin{figure}[htbp]
\centering
\resizebox{\linewidth}{!}{%
\begin{tikzpicture}[
  layer/.style={draw=nublue,fill=nublue!#1,minimum width=10.5cm,minimum height=0.85cm,
    font=\sffamily\small,align=center,rounded corners=2pt},
  node distance=1.5mm]
\node[layer=5]  (l5) {\textbf{Domain-specific languages}: SQL, MATLAB/Simulink, Modelica, spreadsheet formulas};
\node[layer=10, below=of l5] (l4) {\textbf{High-level, interpreted / JIT}: \py, R, JavaScript, Julia, Wolfram};
\node[layer=18, below=of l4] (l3) {\textbf{High-level, compiled}: Fortran, C++, Java, C\#, Rust, Go};
\node[layer=28, below=of l3] (l2) {\textbf{``Portable assembly''}: C};
\node[layer=40, below=of l2] (l1) {\textbf{Assembly language}: \texttt{mov}, \texttt{add}, \texttt{jmp} \ldots};
\node[layer=55, below=of l1,text=white] (l0) {\textbf{Machine code}: \texttt{10110000 01100001} \ldots};
\draw[-{Stealth[length=3mm]},thick,nugold] ([xshift=6mm]l0.east) -- node[right,font=\sffamily\small,align=left]{more\\abstract,\\closer to\\the problem} ([xshift=6mm]l5.east);
\draw[-{Stealth[length=3mm]},thick,nugold] ([xshift=-6mm]l5.west) -- node[left,font=\sffamily\small,align=right]{more\\control,\\closer to\\the hardware} ([xshift=-6mm]l0.west);
\end{tikzpicture}}
\caption{Levels of programming languages. Higher levels let you express the problem
more directly; lower levels give you more control over what the hardware does.}
\label{fig:levels}
\end{figure}

\subsection*{The processor understands only machine code}

To see what ``low level'' means, consider adding two numbers. In \py you write
\code{c = a + b}. In x86-64 assembly language the same thing could look like this:
\begin{lstlisting}[language={[x86masm]Assembler}]
    mov  eax, DWORD PTR [rbp-4]    ; load a from memory into register eax
    add  eax, DWORD PTR [rbp-8]    ; add b to it
    mov  DWORD PTR [rbp-12], eax   ; store the result in c
\end{lstlisting}
Each line is one instruction for the processor. An \term{assembler} translates each
instruction into its binary machine code. Assembly is specific to one processor family:
code for an Intel laptop does not run on the ARM processor in your phone. Today
assembly is used only where every clock cycle or every byte counts, for example in the
start-up code of a microcontroller or in the innermost loops of a numerical library.

\subsection*{A high-level language hides the machine}

A \term{high-level language} lets you write \code{c = a + b} and leaves the details
(which registers to use, where in memory to put \code{c}) to a translator program. The
first widely used high-level language was \textbf{Fortran} (1957), designed at IBM for
scientists and engineers who wanted to write formulas, not machine instructions. Its name
is short for \emph{FORmula TRANslation}. Chapter~\ref{ch:fortran} tells its story.

High-level languages bring three big advantages:
\begin{itemize}
\item \textbf{Productivity}: one line of high-level code can replace dozens of machine
  instructions.
\item \textbf{Portability}: the same source code can be translated for different
  processors.
\item \textbf{Readability}: other people (and you, six months later) can understand the
  program.
\end{itemize}
The price is that you give up some control. The translator decides how to use the
hardware, and it may not make the choice you would have made. For most engineering work
this is a very good trade.

\section{Programming paradigms}

A \term{paradigm} is a style of organising a program. Most modern languages support
several paradigms, but each language has a ``natural'' style. Knowing the paradigms helps
you understand why languages look the way they do.

\subsection*{Imperative programs change state step by step}

In the \term{imperative} style a program is a sequence of commands that change the
\emph{state} of the computer (the values stored in variables). ``Set \code{x} to 0. Add 1
to \code{x}. If \code{x} is less than 10, repeat.'' \term[procedural programming]{Procedural} programming
organises those commands into named \emph{procedures} (also called functions or
subroutines) that can be called again and again. Fortran, C, and the scripts you write in
the first weeks of ENG 101 are procedural.
\begin{lstlisting}[language=Python]
# Procedural style: a function that changes nothing outside itself
def mean(values):
    total = 0.0
    for v in values:
        total = total + v
    return total / len(values)
\end{lstlisting}

\subsection*{Objects bundle data with the functions that use them}

In \term{object-oriented programming} (OOP) data and the functions that work on that data
are bundled together into \term{objects}. A \term{class} is a blueprint for objects. For
example, a class \code{Pipe} could store a length, a diameter and a roughness, and offer
a method \code{pressure\_drop(flow\_rate)}. Each actual pipe in a network is an object (an
\term{instance}) of that class. C++, Java and C\# are built around OOP; \py supports it
fully.
\begin{lstlisting}[language=Python]
class Pipe:
    def __init__(self, length, diameter):
        self.length = length        # m
        self.diameter = diameter    # m

    def area(self):
        return 3.141592653589793 * self.diameter**2 / 4

p = Pipe(length=12.0, diameter=0.05)
print(p.area())                     # 0.0019634954084936207
\end{lstlisting}
Three ideas usually come with OOP:
\begin{itemize}
\item \term{encapsulation}: the inside of an object is hidden; you use it through its
  methods;
\item \term{inheritance}: a new class can extend an existing one (a \code{HeatedPipe} is
  a \code{Pipe} with a wall temperature);
\item \term{polymorphism}: different classes can offer methods with the same name, and
  code that calls \code{area()} works for any of them.
\end{itemize}

\subsection*{Functional programs combine functions without side effects}

In \term{functional programming} the main building block is the mathematical function:
a rule that maps inputs to outputs and has no side effects. Instead of changing variables
in a loop, you combine functions. Languages such as Haskell, Lisp, OCaml and F\# are
designed for this style; the Wolfram Language and R are strongly influenced by it; and
\py, JavaScript and modern C++ borrow many of its ideas.
\begin{lstlisting}[language=Python]
# Functional style: no variable is changed after it is created
values = [3.1, 2.9, 3.0, 3.2]
squares = list(map(lambda v: v**2, values))
mean = sum(values) / len(values)
\end{lstlisting}
Functional code is often shorter and easier to test, because a function with no side
effects always returns the same output for the same input.

\subsection*{Declarative programs state what is wanted, not how}

In \term{declarative} languages you describe \emph{what} you want, not \emph{how} to get
it. The classic example is SQL, the language of databases:
\begin{lstlisting}[language=SQL]
SELECT sample_id, AVG(yield_strength)
FROM tensile_tests
WHERE material = 'Al6061'
GROUP BY sample_id;
\end{lstlisting}
You do not write a loop over the rows; the database decides how to find and average
them. Spreadsheet formulas, regular expressions, HTML and the equation-based modelling
language Modelica (Section~\ref{sec:modelica}) are also declarative.

\subsection*{Array programs operate on whole vectors and matrices}

Engineers and scientists work with vectors and matrices. \term[array programming]{Array programming} lets
you apply an operation to a whole array at once, without writing a loop. If \code{T} is an
array of 1000 temperatures in Celsius, then in MATLAB, \py with NumPy, R, Julia or modern
Fortran you can write
\begin{lstlisting}[language=Python]
T_kelvin = T + 273.15
\end{lstlisting}
and all 1000 values are converted. This style was pioneered by the language APL in the
1960s and became mainstream through MATLAB. It is the natural style for numerical work,
and it is usually much faster than an explicit loop in an interpreted language
(Section~\ref{sec:numpy-speed}).

\begin{keyidea}
Paradigms are tools, not religions. A good engineering program often mixes them: an
object to represent a reactor, array operations to compute its temperature profile, a
pure function for the reaction rate, and a short imperative loop for time stepping.
\end{keyidea}

\section{Types and type systems}

Every value in a program has a \term{type}: an integer, a floating-point number, a
string of characters, a list, a matrix, a date. The type tells the language how the value
is stored and which operations make sense. You can multiply two numbers, but not two
dates. Languages differ a lot in how they handle types, and these differences explain
much of what you see when you look at code in different languages.

\subsection*{Types can be fixed before the run or carried by the values}

In a \term{statically typed} language the type of every variable is known before the
program runs, usually because you declare it:
\begin{lstlisting}[language=C]
double pressure = 101325.0;   /* C: pressure can only ever hold a double */
int    n_steps  = 40;
\end{lstlisting}
If you later try to store a string in \code{pressure}, the compiler refuses to build the
program. C, C++, C\#, Java, Fortran, Rust and Go are statically typed.

In a \term{dynamically typed} language a variable is just a name attached to a value, and
the value carries its type with it:
\begin{lstlisting}[language=Python]
pressure = 101325.0     # pressure refers to a float
pressure = "1 atm"      # now it refers to a string; Python does not complain
\end{lstlisting}
\py, JavaScript, MATLAB, R and the Wolfram Language are dynamically typed. Julia is
dynamically typed but lets you add type annotations, and uses type information to compile
fast code.

Static typing catches many mistakes early and lets the compiler produce fast code.
Dynamic typing makes code shorter and quicker to write and experiment with. Many modern
languages try to have both: \py allows optional \term{type hints},
\begin{lstlisting}[language=Python]
def terminal_velocity(m: float, c: float, g: float = 9.81) -> float:
    return (m * g / c) ** 0.5
\end{lstlisting}
and TypeScript adds static types to JavaScript (Section~\ref{sec:typescript}).

\subsection*{Some languages convert types silently}

A separate question is how strictly a language enforces types when you mix them. In
\py, \code{"3" + 4} is an error: you cannot add a string and an integer. In JavaScript,
\code{"3" + 4} quietly produces the string \code{"34"}, and \code{"3" * 4} produces the
number \code{12}. We say \py is \term{strongly typed} and JavaScript is more
\term{weakly typed}. Weak typing is convenient for web pages but can hide errors in
calculations. C is statically typed but fairly weak: it will silently convert a
\code{double} to an \code{int} and throw away the fractional part.

\begin{pitfall}[Common pitfall: silent conversions]
In C, Java, C\# and older Fortran, dividing two integers gives an integer:
\code{7 / 2} is \code{3}, not \code{3.5}. In \py~3, \code{7 / 2} is \code{3.5}
and \code{7 // 2} is \code{3}. In MATLAB, numbers are floating point by default, so
\code{7/2} is \code{3.5}. When you translate a formula from one language to another,
check every division.
\end{pitfall}

\section{Compiled and interpreted languages}

The source code you write has to be translated into machine code before the processor can
run it. There are two main strategies, and Chapter~\ref{ch:howcoderuns} explains them
in detail.
\begin{itemize}
\item A \term{compiler} translates the whole program into machine code \emph{before} it
  runs, and saves the result as an executable file. You run the executable. Fortran, C,
  C++, Rust and Go are usually compiled.
\item An \term{interpreter} reads the program and carries out its instructions
  \emph{while} it runs. \py (in its standard implementation), R and MATLAB are usually
  interpreted.
\end{itemize}
Between these two extremes are languages that compile to an intermediate
\term{bytecode} that runs on a \term{virtual machine} (Java, C\#), and languages that
compile pieces of the program to machine code just before they are needed, called
\term{just-in-time (JIT) compilation} (Julia, JavaScript in a web browser, MATLAB's
modern engine). Strictly speaking, ``compiled'' and ``interpreted'' describe
\emph{implementations}, not languages: there are \py compilers and C interpreters. But
each language has a usual way of being run, and that usual way shapes how it feels to
use it.

\section{A short history of programming languages}

Figure~\ref{fig:timeline} shows when the languages in these notes first appeared. A few
milestones are worth knowing.

\begin{figure}[htbp]
\centering
\begin{tikzpicture}[x=0.152cm,y=0.5cm,
  lang/.style={font=\sffamily\scriptsize,anchor=west,inner sep=1pt},
  dot/.style={circle,fill=nublue,inner sep=1.4pt}]
\draw[thick,gray] (1950,0) -- (2025,0);
\foreach \y in {1950,1960,...,2020} {
  \draw[gray] (\y,0.15) -- (\y,-0.15) node[below,font=\sffamily\scriptsize]{\y};
}
% name, year, row
\foreach \name/\yr/\row in {
   Fortran/1957/1, LISP/1958/2, COBOL/1959/3, BASIC/1964/4, APL/1966/5,
   Simula/1967/6, C/1972/1, Prolog/1972/2, SQL/1974/3, S (\textrightarrow\ R)/1976/4,
   MATLAB/1984/5, C++/1985/6, Perl/1987/7, Mathematica/1988/8,
   Python/1991/1, Visual Basic/1991/2, R/1993/3, Java/1995/4, JavaScript/1995/5,
   PHP/1995/6, C\#/2000/7, Scala/2004/8, Go/2009/1, Rust/2010/2, Julia/2012/3,
   TypeScript/2012/4, Swift/2014/5, Mojo/2023/6} {
  \draw[gray!45] (\yr,0) -- (\yr,\row);
}
\foreach \name/\yr/\row in {
   Fortran/1957/1, LISP/1958/2, COBOL/1959/3, BASIC/1964/4, APL/1966/5,
   Simula/1967/6, C/1972/1, Prolog/1972/2, SQL/1974/3, S (\textrightarrow\ R)/1976/4,
   MATLAB/1984/5, C++/1985/6, Perl/1987/7, Mathematica/1988/8,
   Python/1991/1, Visual Basic/1991/2, R/1993/3, Java/1995/4, JavaScript/1995/5,
   PHP/1995/6, C\#/2000/7, Scala/2004/8, Go/2009/1, Rust/2010/2, Julia/2012/3,
   TypeScript/2012/4, Swift/2014/5, Mojo/2023/6} {
  \node[dot] at (\yr,\row) {};
  \node[lang,fill=white] at (\yr+0.6,\row) {\name};
}
\end{tikzpicture}
\caption{Year of first appearance of selected programming languages. Dates are the
first public release or first published description.}
\label{fig:timeline}
\end{figure}

\begin{description}
\item[1940s] The first electronic computers (ENIAC, Colossus, the Manchester Baby) were
  programmed by setting switches and plugging cables, and then in machine code.
\item[1950s] Grace Hopper built the first compilers and argued that programs should be
  written in something close to English. John Backus's team at IBM released
  \textbf{Fortran} in 1957 and showed that compiled code could be almost as fast as
  hand-written machine code. \textbf{LISP} (1958) introduced many ideas of functional
  programming, and \textbf{COBOL} (1959) became the language of business data
  processing --- it still runs much of the world's banking.
\item[1960s] \textbf{BASIC} (1964) was designed at Dartmouth College so that
  students in every subject could use a computer. \textbf{Simula} (1967), from Norway,
  introduced classes and objects for simulating systems such as queues and traffic.
  \textbf{APL} introduced array programming.
\item[1970s] Dennis Ritchie created \textbf{C} at Bell Labs (1972) to write the Unix
  operating system. C's combination of efficiency and portability made it the
  foundation of most software that exists today. \textbf{SQL} gave databases a common
  language.
\item[1980s] Cleve Moler wrote the first \textbf{MATLAB} so that his students could use
  the Fortran linear-algebra libraries LINPACK and EISPACK without writing Fortran; it
  became a commercial product in 1984. Bjarne Stroustrup added classes to C to create
  \textbf{C++} (1985). Stephen Wolfram released \textbf{Mathematica} (1988), bringing
  symbolic mathematics to the desktop.
\item[1990s] Guido van Rossum released \textbf{\py} (1991). Ross Ihaka and Robert
  Gentleman began \textbf{R} (1993) at the University of Auckland. Sun Microsystems
  released \textbf{Java} (1995) with the promise ``write once, run anywhere'', and
  Brendan Eich wrote the first version of \textbf{JavaScript} for the Netscape browser in
  about ten days (1995).
\item[2000s] Microsoft released \textbf{C\#} and the .NET platform (2000). Scientific
  \py grew around NumPy (2006), and Google released \textbf{Go} (2009).
\item[2010s] Mozilla started \textbf{Rust} (2010; version 1.0 in 2015), a systems
  language that prevents whole classes of memory errors. \textbf{Julia} (2012) was
  designed at MIT for fast scientific computing. Jupyter notebooks (2014) changed how
  scientists share code. Deep learning, mostly driven from \py, made \py the
  dominant language of data science and AI.
\item[2020s] AI coding assistants based on large language models became part of everyday
  programming. They can write and explain code in every language in these notes, which
  makes \emph{reading} and \emph{checking} code an even more important skill
  (Section~\ref{sec:ai-assistants}).
\end{description}

\section{Language families}

Languages borrow from each other. Recognising family resemblances makes a new language
much easier to read.

\begin{itemize}
\item \textbf{The C family} uses braces \code{\{\}} for blocks, semicolons to end
  statements, \code{//} or \code{/* */} for comments, and \code{for (i = 0; i < n; i++)}
  loops. It includes C, C++, C\#, Java, JavaScript, Go, Rust (partly), PHP and many
  others. If you can read one of them, you can roughly read all of them.
\item \textbf{The Fortran / MATLAB family} uses \code{end} to close blocks, indices that
  start at 1, column-major arrays (Section~\ref{sec:rowcolumn}) and built-in matrix
  operations. Julia's syntax is strongly influenced by MATLAB.
\item \textbf{The scripting family} (\py, Ruby, Perl, R, Lua) is dynamically typed and
  interpreted and favours short, readable code.
\item \textbf{The functional family} (LISP, Scheme, Haskell, ML, OCaml, F\#, Erlang,
  Elixir) treats functions as values and avoids changing state. The Wolfram Language
  belongs here in spirit.
\end{itemize}

Table~\ref{tab:hello} shows the classic first program, which prints the words
``Hello, world!'', in each language covered in these notes.

\begin{table}[htbp]
\centering
\caption{``Hello, world!'' in the languages of these notes. Notice the amount of
\emph{ceremony} (code that is needed but does not do the actual work) in each.}
\label{tab:hello}
\small
\begin{tabularx}{\linewidth}{@{}l>{\raggedright\arraybackslash\ttfamily\footnotesize}X@{}}
\toprule
\textbf{Language} & \normalfont\textbf{Program}\\
\midrule
Python      & print("Hello, world!")\\[3pt]
Fortran     & program hello\newline\hspace*{1em}print *, "Hello, world!"\newline end program hello\\[3pt]
C           & \#include <stdio.h>\newline int main(void) \{\newline\hspace*{1em}printf("Hello, world!\textbackslash n");\newline\hspace*{1em}return 0;\newline\}\\[3pt]
C++         & \#include <iostream>\newline int main() \{\newline\hspace*{1em}std::cout << "Hello, world!\textbackslash n";\newline\}\\[3pt]
C\#         & Console.WriteLine("Hello, world!");\hfill\normalfont\textit{(top-level statement, .NET 6+)}\\[3pt]
Java        & public class Hello \{\newline\hspace*{1em}public static void main(String[] args) \{\newline\hspace*{2em}System.out.println("Hello, world!");\newline\hspace*{1em}\}\newline\}\\[3pt]
JavaScript  & console.log("Hello, world!");\\[3pt]
MATLAB      & disp(\textquotesingle{}Hello, world!\textquotesingle{})\\[3pt]
Wolfram     & Print["Hello, world!"]\\[3pt]
R           & print("Hello, world!")\\[3pt]
Julia       & println("Hello, world!")\\[3pt]
Rust        & fn main() \{\newline\hspace*{1em}println!("Hello, world!");\newline\}\\[3pt]
Go          & package main\newline import "fmt"\newline func main() \{\newline\hspace*{1em}fmt.Println("Hello, world!")\newline\}\\[3pt]
SQL         & SELECT \textquotesingle{}Hello, world!\textquotesingle{};\\[3pt]
Bash        & echo "Hello, world!"\\
\bottomrule
\end{tabularx}
\end{table}

\section{Which engineers use which languages?}

There is no single ``engineering language''. Table~\ref{tab:disciplines} gives a rough
picture of where each language is commonly met. It is a generalisation; every company and
research group has its own habits.

\begin{table}[htbp]
\centering
\caption{Where engineers commonly meet each language (a rough guide).}
\label{tab:disciplines}
\small
\begin{tabularx}{\linewidth}{@{}l>{\raggedright\arraybackslash}X@{}}
\toprule
\textbf{Field} & \textbf{Typical languages and tools}\\
\midrule
Chemical \& process & \py (data, process models), MATLAB, Excel/VBA, Aspen Plus and gPROMS
  (own modelling languages), Fortran (legacy thermodynamics), Modelica\\
Materials science & \py (ASE, pymatgen, data analysis), Fortran (VASP, Quantum ESPRESSO),
  C++ (LAMMPS), Julia, MATLAB\\
Mechanical \& aerospace & MATLAB/Simulink, \py, C/C++ (embedded, simulation), Fortran (CFD,
  FEM), Modelica\\
Civil \& structural & \py, Excel/VBA, MATLAB, scripting in FEM packages (\py in Abaqus,
  APDL in ANSYS), Dynamo/C\# in BIM tools\\
Electrical \& electronic & C (microcontrollers), VHDL/Verilog (hardware), MATLAB, \py,
  C++, assembly\\
Robotics \& mechatronics & C++ and \py (ROS), C (firmware), MATLAB, Structured Text (PLCs)\\
Petroleum \& mining & \py, MATLAB, Fortran/C++ (reservoir simulators), R (geostatistics)\\
Biomedical & \py, MATLAB, R, C (devices)\\
Software \& computer & everything; especially C, C++, Java, C\#, JavaScript/TypeScript,
  \py, Go, Rust, SQL\\
Data science \& AI (all fields) & \py (NumPy, pandas, PyTorch), R, SQL, Julia\\
Manufacturing \& automation & Structured Text and Ladder Logic (PLCs), G-code (CNC),
  \py, C\#\\
\bottomrule
\end{tabularx}
\end{table}

\begin{engnote}[Engineering note: the language is not the product]
Clients do not buy code; they buy an answer, a design or a working machine. A few
lines of \py that produce a correct, checked result in an afternoon are often
worth more than a fast C++ program that takes a month. On the other hand, the control
software of an aircraft must be fast, predictable and certified, and there \py is
simply not an option. Choosing a language is part of the engineering design.
\end{engnote}

\section{How popular is each language?}
\label{sec:popularity}

Students often ask: \emph{``Which language is the most popular?''} It is a fair question,
but the answer depends on what you count. This section looks at one well-known ranking, the
TIOBE index, explains how to read it, and shows why different rankings disagree.

\subsection*{The TIOBE index counts mentions on the web}

The \term{TIOBE index} is a list published every month by the Dutch software-quality
company TIOBE (\url{https://www.tiobe.com/tiobe-index/}). It is based on searches of the
web: roughly speaking, it counts how many web pages, courses and suppliers mention each
programming language, using several large search engines and websites. Each language then
receives a \emph{rating}: its share of all those hits, in per cent. A rating of 20\% for
\py therefore means that about one fifth of the hits counted for \emph{all} languages
together were about \py. It does \emph{not} mean that 20\% of all programs are written in
\py, and it says nothing about which language is ``best''.

Table~\ref{tab:tiobe} shows the top 20 for November 2026, the month in which these notes
were written, together with the place in these notes where each language is discussed.

\begin{table}[htbp]
\centering
\caption{The 20 highest-rated languages in the TIOBE index for November 2026, and where
each one appears in these notes. ``Change'' is the change in rating compared with
November 2025, in percentage points.}
\label{tab:tiobe}
\small
\begin{tabular}{@{}rlrrl@{}}
\toprule
\textbf{Rank} & \textbf{Language} & \textbf{Rating} & \textbf{Change} & \textbf{In these notes}\\
\midrule
1 & \py & 20.86\% & $-2.51$ & Chapter~\ref{ch:python}\\
2 & C & 11.48\% & $+1.80$ & Chapter~\ref{ch:c}\\
3 & Java & 8.89\% & $+0.34$ & Chapter~\ref{ch:java}\\
4 & C++ & 8.35\% & $-0.60$ & Chapter~\ref{ch:cpp}\\
5 & C\# & 3.39\% & $-4.26$ & Chapter~\ref{ch:csharp}\\
6 & JavaScript & 2.42\% & $-1.00$ & Chapter~\ref{ch:javascript}\\
7 & R & 2.27\% & $+0.60$ & Chapter~\ref{ch:r}\\
8 & SQL & 1.91\% & $+0.11$ & Chapter~\ref{ch:others}\\
9 & Visual Basic & 1.88\% & $-1.43$ & Chapters~\ref{ch:csharp} and \ref{ch:others} (VBA)\\
10 & Go & 1.69\% & $-0.03$ & Chapter~\ref{ch:others}\\
11 & Rust & 1.63\% & $+0.25$ & Chapter~\ref{ch:others}\\
12 & PHP & 1.58\% & $+0.27$ & Chapter~\ref{ch:others} (briefly)\\
13 & Fortran & 1.21\% & $-0.34$ & Chapter~\ref{ch:fortran}\\
14 & Scratch & 1.15\% & $+0.13$ & Chapter~\ref{ch:others} (briefly)\\
15 & Perl & 1.14\% & $-0.69$ & Chapter~\ref{ch:others} (briefly)\\
16 & MATLAB & 1.08\% & $-0.30$ & Chapter~\ref{ch:matlab}\\
17 & Assembly language & 0.92\% & $-0.24$ & Chapters~\ref{ch:intro} and \ref{ch:others}\\
18 & Swift & 0.89\% & $+0.03$ & Chapter~\ref{ch:others} (briefly)\\
19 & Delphi/Object Pascal & 0.78\% & $-1.28$ & Chapter~\ref{ch:others} (briefly)\\
20 & Dart & 0.74\% & $+0.05$ & Chapter~\ref{ch:others} (briefly)\\
\bottomrule
\end{tabular}
\end{table}

\subsection*{Old and general-purpose languages lead the table}

Four observations are worth making.

\begin{enumerate}
\item \textbf{\py is clearly first.} Its rating is almost twice that of the second
  language. This is one reason why ENG 101 teaches \py: the skills you learn are the ones
  most in demand, and help is easy to find.
\item \textbf{Old languages are still strong.} C (1972) is second, Fortran (1957) is
  13th, and assembly language is 17th. Languages do not disappear when newer ones arrive,
  because the software written in them is still running and still needs engineers who can
  read it.
\item \textbf{The top four take half of everything.} \py, C, Java and C++ together have
  a rating of about 50\%. After the fifth place every language is below 3\%, so the
  differences between, say, rank 10 and rank 20 are small and can change from month to
  month.
\item \textbf{Low rank does not mean unimportant.} Julia is 23rd (0.68\%), Ladder Logic
  is 37th, LabVIEW is 48th, and the Wolfram Language, Structured Text and VHDL are
  somewhere between 51 and 100. Yet each of these is the standard tool in its own field: a
  control engineer programs PLCs in Ladder Logic every day, whatever the index says.
  Specialised languages have few users compared with general-purpose ones, but those
  users cannot replace them.
\end{enumerate}

\subsection*{Other rankings measure something different}

You may read elsewhere that JavaScript and TypeScript are the most popular languages. That
statement is also based on real data --- it just measures something different.
\begin{itemize}
\item \textbf{TIOBE} counts mentions on the web. It favours languages that are widely
  \emph{taught and written about}, including old ones. Here JavaScript is 6th and
  TypeScript only 41st.
\item \textbf{Rankings based on GitHub}, the website where most open-source code is
  stored, count how much code is \emph{being written and shared} in public projects.
  Web development produces an enormous amount of public code, so JavaScript, TypeScript
  and \py come out on top.
\item \textbf{Developer surveys}, such as the yearly Stack Overflow survey, ask
  programmers which languages they \emph{use}. Most respondents are professional software
  developers, many of them web developers, so JavaScript, SQL and \py lead.
\item \textbf{Job advertisements} show what employers in a particular country and
  industry are looking for, which may be different again.
\end{itemize}
None of these rankings counts the code that engineers write most often: scripts inside a
company, firmware in a product, PLC programs in a factory, or simulation codes on a
cluster. This code is rarely published on the web, so it is almost invisible to every
ranking.

\begin{keyidea}
A popularity ranking measures how visible a language is in one particular place (web
pages, public code, a survey). It is useful for answering ``Will I find help, libraries
and colleagues for this language?'' It does not answer ``Which language should I use for
this task?'' For that question, use the engineering criteria of
Chapter~\ref{ch:choosing}.
\end{keyidea}

\begin{tryit}
Open \url{https://www.tiobe.com/tiobe-index/} and compare today's top 10 with
Table~\ref{tab:tiobe}. Which languages have moved? The page also has a chart of ratings
since 2001: find the year in which \py overtook Java and C.
\end{tryit}

\section{What all languages share}

Despite their differences, every general-purpose language provides the same core building
blocks. Chapter by chapter we will look for them:
\begin{enumerate}
\item \textbf{values and types} --- numbers, text, true/false;
\item \textbf{variables} --- names that refer to values;
\item \textbf{operators and expressions} --- \code{+}, \code{*}, \code{<}, \code{and}
  \ldots;
\item \textbf{control flow} --- \code{if}/\code{else} decisions and loops;
\item \textbf{functions} --- named, reusable pieces of code with inputs and outputs;
\item \textbf{collections} --- arrays, lists, dictionaries;
\item \textbf{input and output} --- reading and writing the screen and files;
\item \textbf{modules and libraries} --- code written by others that you can reuse;
\item \textbf{error handling} --- what happens when something goes wrong.
\end{enumerate}
Each language chapter is organised around this list, so you can jump straight to, say,
``functions in Fortran'' and compare it with ``functions in \py''.

\begin{exercises}
\item Write the Euclidean algorithm for the greatest common divisor as a flowchart. Then
  trace it by hand for $a = 252$, $b = 105$. How many times does the loop run?
\item Classify each of the following as imperative, object-oriented, functional,
  declarative or array style, and explain why: (a) an Excel formula
  \code{=AVERAGE(B2:B50)}; (b) a \py loop that adds numbers one by one; (c) the SQL query in
  this chapter; (d) the MATLAB statement \code{y = A*x + b}; (e) a \py class
  \code{Reactor} with a method \code{step(dt)}.
\item The text says that ``compiled'' and ``interpreted'' describe implementations, not
  languages. Find one example of a compiler for \py and one example of an interpreter for C
  (a web search will help). Why do you think they are rarely used?
\item Look at Table~\ref{tab:hello}. Count the lines needed for ``Hello, world!'' in each
  language. Which languages need a function called \code{main}? Why do you think compiled
  languages tend to need one?
\item Ask an engineer you know (a lecturer, a relative, an alumnus) which programming
  languages they use at work, and for what. Compare the answer with
  Table~\ref{tab:disciplines}.
\item In JavaScript, \code{"3" + 4} gives \code{"34"}, but \code{"3" - 4} gives \code{-1}.
  Explain in your own words why such behaviour could be dangerous in an engineering
  calculation.
\item Look at Table~\ref{tab:tiobe}. Add up the ratings of the languages that have their
  own chapter in these notes. What share of the TIOBE index do these chapters cover? Then
  explain in two or three sentences why TypeScript can be near the top of a ranking based on
  GitHub but only 41st in the TIOBE index.
\end{exercises}
